\documentclass[11pt,letterpaper]{article}

\usepackage[top=1in, bottom=1in, left=1in, right=1in, headheight=14pt]{geometry}
\usepackage{fontspec}
\setmainfont{DejaVu Serif}
\setsansfont{DejaVu Sans}
\setmonofont{DejaVu Sans Mono}

\usepackage{xcolor}
\usepackage{titlesec}
\usepackage{fancyhdr}
\usepackage{enumitem}
\usepackage{hyperref}
\usepackage{booktabs}
\usepackage{tabularx}
\usepackage{longtable}
\usepackage{colortbl}
\usepackage{multirow}
\usepackage{array}
\usepackage{parskip}
\usepackage[breakable,skins,listings]{tcolorbox}
\usepackage{listings}
\usepackage{graphicx}
\usepackage{lastpage}

% ── Color Scheme: Broadcast Heritage (History domain) ─────────────────────────
\definecolor{primary}{HTML}{B71C1C}       % Deep broadcast red
\definecolor{secondary}{HTML}{E65100}     % Vacuum-tube amber
\definecolor{tertiary}{HTML}{1A237E}      % FCC regulatory navy
\definecolor{accent}{HTML}{C62828}        % Active accent
\definecolor{lightprimary}{HTML}{FFEBEE}  % Red tint bg
\definecolor{lightsecondary}{HTML}{FFF3E0}
\definecolor{lighttertiary}{HTML}{E8EAF6}
\definecolor{rowgray}{HTML}{F5F5F5}
\definecolor{bordergray}{HTML}{BDBDBD}
\definecolor{subtlegray}{HTML}{757575}
\definecolor{darktext}{HTML}{212121}

% ── Section Formatting ─────────────────────────────────────────────────────────
\titleformat{\section}
  {\Large\bfseries\sffamily\color{primary}}
  {\thesection}{1em}{}
  [\vspace{-0.5em}\textcolor{bordergray}{\rule{\textwidth}{0.4pt}}]

\titleformat{\subsection}
  {\large\bfseries\sffamily\color{secondary}}
  {\thesubsection}{1em}{}

\titleformat{\subsubsection}
  {\normalsize\bfseries\sffamily\color{tertiary}}
  {\thesubsubsection}{1em}{}

\titlespacing*{\section}{0pt}{1.5em}{0.8em}
\titlespacing*{\subsection}{0pt}{1.2em}{0.5em}
\titlespacing*{\subsubsection}{0pt}{0.8em}{0.3em}

% ── Custom Environments ────────────────────────────────────────────────────────
\newcommand{\stageheader}[2]{%
  \noindent\colorbox{#2}{\makebox[\textwidth][l]{%
    \textcolor{white}{\bfseries\sffamily\hspace{6pt}#1\hspace{6pt}}}}%
  \vspace{0.5em}\par%
}

\lstdefinestyle{asciiStyle}{
  basicstyle=\small\ttfamily,
  breaklines=true,
  breakatwhitespace=false,
  columns=fullflexible,
  keepspaces=true,
}
\newtcblisting{asciibox}{
  colback=rowgray, colframe=bordergray,
  arc=1pt, boxrule=0.5pt,
  left=6pt, right=6pt, top=4pt, bottom=4pt,
  breakable,
  listing only,
  listing style=asciiStyle,
}

\newtcolorbox{quotebox}{
  colback=lightprimary, colframe=primary,
  arc=2pt, boxrule=0.5pt,
  left=12pt, right=12pt, top=8pt, bottom=8pt,
  breakable
}

\newcommand{\metafield}[2]{\textbf{\sffamily #1} #2\par}

% ── Table Helpers ──────────────────────────────────────────────────────────────
\newcolumntype{L}[1]{>{\raggedright\arraybackslash}p{#1}}
\newcolumntype{C}[1]{>{\centering\arraybackslash}p{#1}}
\newcommand{\tableheader}[1]{\textbf{\sffamily\textcolor{white}{#1}}}

% ── Headers & Footers ──────────────────────────────────────────────────────────
\pagestyle{fancy}
\fancyhf{}
\fancyhead[L]{\small\sffamily\textcolor{subtlegray}{AM Radio: West Coast \& PNW Deep History}}
\fancyhead[R]{\small\sffamily\textcolor{subtlegray}{GSD Mission Package}}
\fancyfoot[C]{\small\sffamily\textcolor{subtlegray}{Page \thepage\ of \pageref{LastPage}}}
\renewcommand{\headrulewidth}{0.4pt}
\renewcommand{\footrulewidth}{0pt}

% ── Hyperlinks ────────────────────────────────────────────────────────────────
\hypersetup{
  colorlinks=true,
  linkcolor=secondary,
  urlcolor=secondary,
  pdfauthor={GSD Ecosystem / Fox Infrastructure Group},
  pdftitle={AM Radio -- West Coast \& PNW Deep Research Mission Package},
  pdfsubject={Vision to Mission Pipeline},
}

% ──────────────────────────────────────────────────────────────────────────────
\begin{document}

% ── TITLE PAGE ────────────────────────────────────────────────────────────────
\thispagestyle{empty}
\begin{center}
\vspace*{3cm}

{\small\sffamily\textcolor{primary}{\textbf{GSD RESEARCH MISSION PACKAGE}}}

\vspace{0.3cm}
\textcolor{bordergray}{\rule{8cm}{0.4pt}}

\vspace{1.5cm}
{\fontsize{28}{34}\selectfont\bfseries\sffamily\textcolor{primary}{The AM Radio Dial\\[0.3em]A Broadcast Heritage Study}}

\vspace{0.8cm}
{\Large\sffamily\textcolor{secondary}{West Coast · Pacific Northwest · California\\History, Licensing, Spectrum \& Future}}

\vspace{0.5cm}
{\large\sffamily\itshape\textcolor{tertiary}{A Deep Research Study}}

\vspace{3cm}
{\sffamily\textcolor{accent}{Vision $\rightarrow$ Research $\rightarrow$ Mission}}\\[0.3em]
{\small\sffamily\textcolor{accent}{Full Pipeline Execution}}

\vspace{3cm}
{\small\sffamily\textcolor{subtlegray}{Date: March 27, 2026  |  Status: Ready for GSD Orchestrator}}\\[0.3em]
{\small\sffamily\textcolor{subtlegray}{Activation Profile: Fleet  |  Estimated: 18--24 context windows across 4--6 sessions}}
\end{center}
\newpage

% ── TABLE OF CONTENTS ─────────────────────────────────────────────────────────
\tableofcontents
\newpage

% ══════════════════════════════════════════════════════════════════════════════
% STAGE 1 — VISION DOCUMENT
% ══════════════════════════════════════════════════════════════════════════════
\stageheader{STAGE 1 --- VISION DOCUMENT}{primary}

\section{AM Radio: West Coast \& PNW --- Vision Guide}

\metafield{Date:}{March 27, 2026}
\metafield{Status:}{Initial Vision / Pre-Research Phase}
\metafield{Depends On:}{None (foundational research mission)}
\metafield{Context:}{Fox Infrastructure Group / GSD Ecosystem Research Pipeline}

\subsection{Vision}

Before the internet, before satellite, before the smartphone---there was a voice in the darkness,
carrying itself over hundreds of miles of ocean, mountain, and forest on invisible waves.
AM radio was not merely a technology; it was the central nervous system of the American West.
It connected loggers in the Cascades to Portland markets, fishermen off the Washington coast to
Seattle weather reports, and ranch hands in the Central Valley to the baseball games of the Giants.
The Pacific Northwest and California were, from the earliest days of licensed broadcasting, a
laboratory for the entire medium.

This research mission traces that arc from the first licensed experimental transmissions in 1920
San Jose to the all-digital AM frontier of 2026. It maps the ecosystem of the AM dial as it
developed along the West Coast---the network affiliations, the format wars, the clear-channel
titans and the daytimer underdogs---with particular focus on the stations of Washington, Oregon,
and California that shaped an entire region's sonic identity. The same curiosity that asks
``what was broadcasting before'' also asks ``what is broadcasting becoming,'' and this mission
addresses both with equal rigor.

The practical thread runs alongside the historical: the mechanics of obtaining a broadcast
license today, the open spectrum that remains available, the rules and regulations under which
all broadcasters must operate, and the pathways---full-power AM, full-power FM, Low Power FM,
FM translators, and Part~15 unlicensed devices---that a new voice can use to reach the air in
2026. The AM dial is not a closed museum; it is a living infrastructure with legal on-ramps
still accessible to those who understand the system.

For the builder who hears frequencies as possibility, this document is a flight manual.

\subsection{Problem Statement}

\begin{enumerate}[leftmargin=*, label=\textbf{\arabic*.}]
  \item \textbf{Fragmented institutional memory.} The history of West Coast AM radio exists in
        scattered archives, fan sites, local historical societies, and FCC dockets. No single
        authoritative synthesis covers the full arc from experimental licenses in 1920 to the
        digital transition of the 2020s with geographic specificity to Washington and California.

  \item \textbf{Licensing opacity for prospective broadcasters.} The FCC licensing process for
        AM, FM, LPFM, and FM translator stations involves multiple forms, filing windows, engineering
        studies, and regulatory conditions that are individually documented but difficult to navigate
        as an integrated pathway. The spectrum scarcity situation, particularly in urban West Coast
        markets, is poorly understood by most non-industry applicants.

  \item \textbf{Unclear open-spectrum inventory.} As of 2026, the actual availability of AM
        frequencies in specific PNW and California markets---accounting for Class A clear-channel
        protections, daytimer restrictions, and NARBA international agreements---is unknown to
        most prospective licensees and not synthesized in accessible form.

  \item \textbf{Regulatory landscape complexity.} Broadcast regulations under 47 CFR Part~73
        span content obligations (obscenity, indecency, EAS), technical standards (power limits,
        antenna systems, modulation), ownership rules (local market caps, cross-ownership), and
        public interest requirements---a body of law that has changed substantially since the
        Telecommunications Act of 1996 and continues to evolve under AM Revitalization proceedings.

  \item \textbf{Uncertainty about AM's future viability.} The EV interference problem, declining
        automotive AM reception, streaming competition, and the HD Radio digital transition create
        genuine uncertainty about the investment value of AM broadcasting infrastructure in 2026
        and beyond---uncertainty that requires evidence-based analysis, not speculation.
\end{enumerate}

\subsection{Core Concept}

\textbf{The Broadcast Archaeology Model:} Treat the AM dial as a stratigraphic record---each
layer of regulation, technology, and format revealing the social and technical forces that shaped
it. Map both the deep past and the navigable present on the same coordinate system.

\subsubsection{West Coast Broadcast Geography}

\begin{asciibox}
WEST COAST AM BROADCAST GEOGRAPHY
══════════════════════════════════════════════════════════

  PACIFIC NORTHWEST ZONE                 CALIFORNIA ZONE
  ┌──────────────────────┐              ┌────────────────────────┐
  │  SEATTLE             │              │  SAN FRANCISCO BAY     │
  │  KJR (950) est.1921  │              │  KGO/KSFO (810) 1924   │
  │  KNWN/KOMO (1000)    │              │  KNBR (680) 1922       │
  │  KIRO (710)          │              │  KYA (1260) 1926       │
  │  KVI (570)           │              │  KCBS (740) 1909*      │
  │  KLAY/KOL            │              │                        │
  │                      │              │  LOS ANGELES           │
  │  TACOMA              │              │  KFI (640) 1922        │
  │  KMO  KTCL           │              │  KNX (1070) 1920       │
  │                      │              │  KABC (790) 1925       │
  │  PORTLAND            │              │  KMPC (710) 1922       │
  │  KGW (620) 1922      │              │                        │
  │  KOIN (940)          │              │  SAN DIEGO             │
  │  KEX (1190)          │              │  KFSD/KOGO (600)       │
  │                      │              │  KCBQ (1170)           │
  │  SPOKANE             │              └────────────────────────┘
  │  KHQ  KGA  KXLY      │
  └──────────────────────┘

  FREQUENCY CLASSES (FCC AM)
  Class A (Clear Channel) — 50 kW, nationwide nighttime reach
  Class B (Regional)      — Up to 50 kW, regional protection
  Class C (Local)         — Up to 1 kW
  Class D (Daytimer)      — Daytime only, sunset-to-sunrise silent

  * KCBS traces origins to experimental station FN in 1909
\end{asciibox}

\subsubsection{Six-Module Research Architecture}

\begin{asciibox}
RESEARCH MODULE ARCHITECTURE
══════════════════════════════════════════════════════════════

  M1: AM BAND          M2: PNW HISTORY       M3: CALIFORNIA
  FOUNDATIONS     ──►  STATIONS         ──►  STATIONS
  (Regulation,         (1920s–present,        (1920s–present,
  Technology,          WA/OR/ID focus)        Bay Area + LA)
  Band Structure)
       │                    │                     │
       ▼                    ▼                     ▼
  M4: SPECTRUM &      M5: FCC RULES &       M6: DIGITAL
  LICENSING      ──►  REGULATIONS      ──►  FUTURE
  PATHWAYS             (47 CFR Pt 73,         (HD Radio,
  (AM/FM/LPFM,         Ownership, EAS,        AM Revival,
  Translators)         Content)               Streaming)

  DEPENDENCY CHAIN: M1 → M2 → M3 → M4 → M5 → M6
  PARALLEL TRACKS:  M2 ‖ M3,  M4 ‖ M5
\end{asciibox}

\subsection{Research Modules}

\subsubsection{Module 1 --- AM Band Foundations}
The physics of amplitude modulation, the history of the AM band (535--1705~kHz), the Federal
Radio Commission of 1927 and its successor the FCC, the Radio Act of 1927, the Communications
Act of 1934, the NARBA international frequency assignment treaty, and the AM station classification
system (Class A through D). This module establishes the technical and regulatory substrate on
which all other modules rest.

\subsubsection{Module 2 --- Pacific Northwest Historical Stations}
A comprehensive station-by-station chronicle of AM broadcasting in Washington, Oregon, and Idaho
from the earliest experimental licenses (KJR 1921, KGW Portland 1922) through the Golden Age of
network radio, the format revolution of Top~40, the news-talk pivot, and the consolidation era
following the Telecommunications Act of 1996. Key stations: KJR, KOMO, KIRO, KVI, KGW, KOIN,
KEX, KHQ, KGA, KOL, and regional stations in Tacoma, Bellingham, Everett, Spokane, and Eugene.

\subsubsection{Module 3 --- California Historical Stations}
The parallel chronicle for California, with emphasis on the Bay Area (KGO/KSFO, KNBR, KYA, KCBS)
and Los Angeles (KFI, KNX, KABC, KMPC, KFWB). The Don Lee Network, the NBC Pacific Coast
network, the NBC Orange Chain, and their West Coast infrastructure. California-specific AM
demographics: as of the AM Revitalization NPRM, 34\% of California AM stations serve specific
minority communities, including 33 Spanish-language, 13 Asian-format, and 4 Korean-format outlets.

\subsubsection{Module 4 --- Spectrum Availability and Licensing Pathways}
The practical guide to entering broadcasting in 2026. Available pathways: (1)~Full-power AM
station---Class A through D, minimum 250~W, requires construction permit; (2)~Full-power FM
station---minimum 100~W, requires open allotment and engineering study; (3)~Low Power FM
(LPFM)---up to 100~W, nonprofit/educational only, 3.5-mile radius, filed during FCC windows;
(4)~FM Translator---rebroadcasts AM or other content on FM; (5)~Part~15 unlicensed---200-foot
radius only. Current spectrum scarcity in Seattle, Portland, San Francisco, and Los Angeles
makes new full-power stations difficult but not impossible.

\subsubsection{Module 5 --- FCC Rules, Regulations, and Broadcaster Obligations}
The regulatory framework governing all licensees under 47 CFR Part~73. Topics: public interest
obligations, the FCC's eight-year license term, Equal Employment Opportunity requirements,
political broadcasting obligations (equal time, lowest unit rate), Emergency Alert System (EAS)
participation, indecency and obscenity restrictions under 18~U.S.C.~\S~1464, ownership caps
(Telecommunications Act of 1996), construction permit requirements, and operational
logs/record-keeping. Recent 2025 FCC rulemaking updating Part~73 references and LPFM separation
requirements.

\subsubsection{Module 6 --- The Digital Future of AM}
The trajectory of AM broadcasting: HD Radio (IBOC hybrid digital), the FCC's 2020 approval of
all-digital AM in the MA3 mode (FCC 20-154), the AM Revitalization proceeding (Docket 13-249),
FM translator access granted to AM stations, the EV radio interference controversy (2023--2024),
and the competitive landscape with streaming audio (AM/FM share fell to 36\% of audio listening
by 2023 per Edison Research). The 95 million vehicles already equipped with HD Radio technology
as a counterweight.

\subsection{Chipset Configuration}

\begin{asciibox}
chipset:
  name: am_radio_west_coast_fleet
  version: "1.0.0"
  profile: fleet
  modules:
    - id: M1
      name: am_band_foundations
      model: sonnet
      parallel_group: null
      depends_on: []
    - id: M2
      name: pnw_historical_stations
      model: sonnet
      parallel_group: alpha
      depends_on: [M1]
    - id: M3
      name: california_historical_stations
      model: sonnet
      parallel_group: alpha
      depends_on: [M1]
    - id: M4
      name: spectrum_and_licensing
      model: sonnet
      parallel_group: beta
      depends_on: [M1]
    - id: M5
      name: fcc_rules_and_regulations
      model: sonnet
      parallel_group: beta
      depends_on: [M1]
    - id: M6
      name: digital_future_am
      model: opus
      parallel_group: null
      depends_on: [M2, M3, M4, M5]
  synthesis:
    model: opus
    depends_on: [M2, M3, M4, M5, M6]
  publication:
    model: sonnet
    depends_on: [synthesis]
  token_budget_k: 280
  parallel_efficiency: 0.72
</asciibox>

\subsection{Scope Boundaries}

\subsubsection{In Scope (v1.0)}
\begin{itemize}[leftmargin=*]
  \item AM band history on the West Coast (Washington, Oregon, Idaho, California) from 1920 to 2026
  \item Station histories for all Class~A and major Class~B AM stations in the PNW and California
  \item Complete FCC licensing pathways for AM, FM, LPFM, and FM translator stations
  \item Spectrum availability analysis for West Coast markets with emphasis on Seattle, Portland, San Francisco, and Los Angeles metro areas
  \item 47 CFR Part~73 rules applicable to new and existing broadcast licensees
  \item AM digital transition (HD Radio, all-digital MA3) and future viability assessment
  \item Open band frequencies and regulatory context for prospective broadcasters
\end{itemize}

\subsubsection{Out of Scope}
\begin{itemize}[leftmargin=*]
  \item FM station history (except where directly relevant to AM simulcast or translator context)
  \item Television broadcasting history
  \item Stations east of the Rockies except as context for clear-channel frequency assignments
  \item Amateur (ham) radio licensing and operations
  \item International broadcasting or shortwave operations
  \item Satellite radio (SiriusXM) except as competitive context
  \item Podcast and streaming platform analysis beyond AM competitive landscape framing
\end{itemize}

\subsection{Success Criteria}

\begin{enumerate}[leftmargin=*, label=\textbf{SC-\arabic{enumi}.}]
  \item All major West Coast AM stations (Class~A and Class~B, PNW and California) individually documented with founding date, founding licensee, call sign history, frequency assignment history, network affiliations, and format evolution.
  \item AM band structure fully explained: frequency range, channel classification (clear/regional/local), station classes (A through D), power limits, and daytime vs.\ nighttime propagation physics.
  \item FCC licensing pathways for AM, FM, LPFM, FM translator, and Part~15 operation described with step-by-step process, required forms, filing fees, and eligibility requirements.
  \item Spectrum scarcity context for Seattle, Portland, San Francisco, and Los Angeles markets quantitatively documented with reference to current FCC allotment tables.
  \item 47 CFR Part~73 regulatory obligations summarized for new licensees: EAS, indecency rules, political broadcasting obligations, EEO requirements, and public inspection file requirements.
  \item Ownership restriction framework documented: local market numerical caps, cross-ownership rules, and changes since the Telecommunications Act of 1996.
  \item AM digital transition options (HD Radio hybrid, all-digital MA3, FM translator) described with current adoption statistics and regulatory status.
  \item EV radio interference problem documented with technical basis and industry response (NAB, automakers, congressional action through the AM Radio for Every Vehicle Act).
  \item Don Lee Network and NBC Pacific Coast network organizational structure and West Coast station affiliations mapped with dates.
  \item Historical format evolution documented per era: network radio (1920s--1940s), Top~40 (1950s--60s), news-talk pivot (1970s--80s), consolidation era (1990s--present).
  \item All numerical claims (wattage, frequencies, station counts, audience shares) attributed to specific FCC documents, FCC databases, NAB, or peer-reviewed sources.
  \item Practical broadcaster entry guide synthesizes licensing, spectrum, and regulatory information into an actionable road map for a prospective West Coast broadcaster in 2026.
\end{enumerate}

\subsection{Relationship to Other GSD Documents}

\begin{center}
\rowcolors{2}{rowgray}{white}
\begin{tabularx}{\textwidth}{L{3.5cm}L{4cm}X}
\toprule
\rowcolor{primary}
\tableheader{Document} & \tableheader{Relationship} & \tableheader{Notes} \\
\midrule
FCC Part 73 Regulatory Reference & Informational dependency & Primary source for Modules 4 and 5 \\
Fox Infrastructure Group SPC Formation & Business context & Broadcasting license applicant entity context \\
GSD Ecosystem Core & Platform & Mission executed via gsd-skill-creator framework \\
\bottomrule
\end{tabularx}
\end{center}

\subsection{The Through-Line}

There is a reason the word ``broadcast'' means both a way of transmitting radio signals and the
act of scattering seed widely across a field. AM radio was the original broadcast medium---a
single transmitter flinging its signal across thousands of square miles, trusting that whoever
needed it would find it. It required nothing of its listeners but a receiver and attention.

The West Coast, always a place where the frontier impulse found new expression, took to radio
with particular intensity. The Pacific Northwest's geography---deep fjords, mountain ranges,
vast distances between settlements---made radio not a luxury but a necessity. The same topography
that makes the Cascades beautiful makes line-of-sight communication impossible. AM's ability to
follow the curve of the earth, to refract off the ionosphere at night and reach Nome or Nome
reach Seattle, was not just technically useful; it was democratically profound.

The builder who picks up this document in 2026 inherits that infrastructure. The frequencies
are still there, allocated and waiting. The regulations are complex but navigable. The question
is not whether the dial has room for new voices---it always has---but whether those voices know
how to find their way onto it.

\newpage

% ══════════════════════════════════════════════════════════════════════════════
% STAGE 2 — RESEARCH REFERENCE
% ══════════════════════════════════════════════════════════════════════════════
\stageheader{STAGE 2 --- RESEARCH REFERENCE}{secondary}

\section{AM Radio: West Coast \& PNW --- Research Reference}

\subsection{How to Use This Document}

This reference compiles primary and secondary research across all six modules. Each section
provides substantive findings with source attribution. Agents should treat this as the ground
truth for the module they are producing; supplement with additional research only where
specifically noted as incomplete.

\textbf{Key source organizations:}
\begin{itemize}[leftmargin=*]
  \item Federal Communications Commission (FCC) --- fcc.gov (primary regulatory authority)
  \item eCFR / Code of Federal Regulations --- Title 47, Part 73
  \item California Historical Radio Society --- californiahistoricalradio.com
  \item HistoryLink.org --- Washington State history archive
  \item Bay Area Radio Museum \& Hall of Fame --- bayarearadio.org
  \item National Association of Broadcasters (NAB)
  \item Prometheus Radio Project / Common Frequency (LPFM advocacy)
  \item Radio Survivor --- radiosurvivor.com
  \item World Radio History Archive --- worldradiohistory.com
\end{itemize}

% ── MODULE 1: AM BAND FOUNDATIONS ─────────────────────────────────────────────
\subsection{Module 1 Reference: AM Band Foundations}

\subsubsection{The Physics of AM and Band Boundaries}

The AM broadcast band occupies 535 to 1,705~kHz (formerly 535--1,600~kHz, expanded in the 1990s
after years of international negotiation). AM channels are spaced 10~kHz apart, yielding
approximately 117 usable channels. The band was expanded upward to 1,700~kHz by the FCC in the
1990s, with frequencies above 1,600~kHz (the ``expanded band'') reserved for existing stations
that caused significant interference lower in the band.

AM propagation differs fundamentally between day and night. During daylight hours, the signal
travels primarily by ground wave, with a reliable range of up to 200 miles depending on power.
At night, the ionosphere's D layer dissipates, allowing the signal to reflect off the E and F
layers and skip hundreds to thousands of miles. This ``skywave'' propagation is why nighttime AM
can be heard across continents---and why stations sharing a frequency must coordinate power and
antenna directionality to prevent interference.

\subsubsection{Regulatory History: FRC, FCC, and the Radio Acts}

Before 1927, the AM dial was chaos. Hundreds of stations operated at overlapping frequencies
with no mandatory coordination. The Federal Radio Commission (FRC) was established by the Radio
Act of 1927, which for the first time required stations to operate in the public interest,
convenience, and necessity. The FRC implemented General Order 40 in November 1928, introducing
the ``clear channel'' concept and reassigning stations to rationalized frequencies.

The Communications Act of 1934 replaced the FRC with the Federal Communications Commission
(FCC), which has regulated broadcast spectrum ever since. The North American Regional Broadcasting
Agreement (NARBA) of 1941 coordinated AM frequency assignments between the United States, Canada,
and Mexico, shifting many West Coast stations to new dial positions (e.g., KOMO shifted from
950 to 1000~kHz under NARBA).

\subsubsection{AM Station Classification System}

\begin{center}
\rowcolors{2}{rowgray}{white}
\begin{tabularx}{\textwidth}{L{2cm}L{3cm}XL{3.5cm}}
\toprule
\rowcolor{primary}
\tableheader{Class} & \tableheader{Old Class} & \tableheader{Power / Coverage} & \tableheader{Key Protection} \\
\midrule
Class A & Class I-A / I-B & Up to 50,000 W; nationwide nighttime & Highest; clear-channel domination \\
Class B & Class II & Up to 50,000 W; regional & Protected from Class A interference \\
Class C & Class III & Up to 1,000 W; local & Local market only \\
Class D & Class IV & Daytime only; up to 1,000 W at night conditionally & Sunset-to-sunrise silent or power reduction \\
\bottomrule
\end{tabularx}
\end{center}

Class~A stations dominate the AM dial at night. The FCC currently estimates 4,400 licensed
commercial AM stations in the United States as of September 2024 (Federal Register, March 2025).

% ── MODULE 2: PNW HISTORICAL STATIONS ─────────────────────────────────────────
\subsection{Module 2 Reference: Pacific Northwest Historical Stations}

\subsubsection{The First Wave: Experimental Stations, 1920--1925}

Radio came to the Pacific Northwest almost simultaneously with its commercial emergence nationwide.
The first convention of the Northwestern Radio Association was held in Portland on June 25, 1921,
attended by 125 members from amateur radio clubs across Washington, Oregon, Idaho, and Northern
California---at a time when only KDKA in Pittsburgh held a commercial license. Portland's Charles
L.~Austin (call sign 7XF) had been broadcasting phonograph music to ships in the Portland Harbor
from his home laboratory on Mount Tabor, and West Coast inventors were already frustrated that
eastern interests received more regulatory recognition.

KJR Seattle, begun by amateur radio operator Vincent~I.~Kraft, was the first radio station
licensed in the Pacific Northwest. It went on the air as an experimental station in 1921 and was
one of the first commercial licenses in the region. In October 1926, Kraft co-founded the Pacific
Broadcasting Corporation with Frederick~C.~Clift of San Francisco, which created KYA in San
Francisco using the call letters and a new transmitter---a direct business connection between
Seattle and Bay Area broadcasting from the earliest days.

KGB in Tacoma went on the air on April 4, 1922, operated by the Tacoma Daily Ledger and the
William~A.~Mullins Electric Company. Notably, its radio editor was Winifred Dow, one of the few
women in the Pacific Northwest to hold an amateur radio license in that era.

\subsubsection{Key PNW AM Stations: Profiles}

\begin{center}
\rowcolors{2}{rowgray}{white}
\begin{tabularx}{\textwidth}{L{2.2cm}C{1.5cm}L{2.5cm}X}
\toprule
\rowcolor{primary}
\tableheader{Station} & \tableheader{Freq.} & \tableheader{Founded} & \tableheader{Key History} \\
\midrule
KJR Seattle & 950 & 1921 & First PNW license; NBC Blue; ESPN Radio; Classic Rock 1988 \\
KNWN/KOMO & 1000 & 1926 & 50 kW Class A; NBC; Fisher family; Mariners; now Lotus Communications \\
KIRO Seattle & 710 & 1927 & CBS affiliate; all-news; Bonneville; 50 kW Class B \\
KVI Seattle & 570 & 1926 & Sports; personality radio; news 1980; Mariners pre-KOMO \\
KGW Portland & 620 & 1922 & NBC affiliate; Orange Chain; first PNW TV 1952 \\
KOIN Portland & 940 & 1928 & Don Lee/CBS affiliate; longtime Portland institution \\
KEX Portland & 1190 & 1927 & NBC Gold Network; Pacific Gold Chain station \\
KHQ Spokane & 590 & 1922 & Spokane dominant; NBC; one of earliest 50 kW outlets \\
KGA Spokane & 1510 & 1927 & NBC Gold; regional powerhouse \\
KFAE Pullman & --- & 1922 & Educational; Washington State College; first edu. station in WA \\
\bottomrule
\end{tabularx}
\end{center}

\subsubsection{The Adolph Linden Empire and Early Consolidation}

Adolph Linden and Edmund Campbell, directors of Puget Sound Savings \& Loan, founded Northwest
Radio Service in 1926 and acquired KJR from Vincent Kraft. Linden then created the American
Broadcasting Company (ABC)---not the later national network---which by 1929 included stations
in Portland, Spokane, San Francisco, and ``a half-dozen others from here to Chicago.''

This early consolidation foreshadowed the industry's later corporate concentration, and Linden's
story ended badly: he was convicted of embezzlement from Puget Sound Savings \& Loan in 1932
and served five years in the Walla Walla state penitentiary.

\subsubsection{NBC Comes to Seattle: The Network Era, 1926--1950}

In 1926, special AT\&T wires reached Seattle from New York, and KOMO became an affiliate of NBC,
the National Broadcasting Company. Local listeners could now hear the same programs---and the
same national commercials---as people across the country. NBC organized its West Coast operations
into the ``Pacific Coast Network'' with Red and Blue designations. In 1931, NBC reorganized into
the Pacific Gold Network (KEX Portland; KJR Seattle) and the Pacific Orange Network (KFI Los
Angeles; KGO Oakland; KGW Portland; KOMO Seattle; KHQ Spokane).

On September 26, 1937, the Mutual Broadcasting System made its Northwest debut with a 90-minute
inaugural program. The Don Lee Network, which produced programming for Mutual's West Coast
affiliates, included: KOL Seattle, KALE Portland (KPOJ), KMO Tacoma, KORE Eugene, KSLM Salem,
KIT Yakima, KVOS Bellingham, KGY Olympia, KIEM Eureka, KPQ Wenatchee, KRNR Roseburg, and KXRO
Aberdeen---described at the time as the largest regional network in the United States.

\subsubsection{Format Evolution: Top 40, News-Talk, and the FM Defection}

Top~40 radio arrived in the Pacific Northwest in the 1950s, and KJR Seattle became one of the
most dominant personality stations in the region through the 1960s and 1970s. KVI (570~AM)
carried the pre-Mariners Seattle Rainiers baseball and was a powerhouse personality station.
When music fans migrated to FM's superior sound in the 1970s, AM pivoted to news and talk.
KIRO became an all-news powerhouse; KOMO followed with news-talk by 1995.

KJR moved to Classic Rock in 1988, then Sports/ESPN in 1992---a format it held for decades.
KOMO's call letters migrated to KNWN (Northwest News Radio) in February 2022 after Lotus
Communications acquired the Fisher/Sinclair radio properties.

% ── MODULE 3: CALIFORNIA HISTORICAL STATIONS ──────────────────────────────────
\subsection{Module 3 Reference: California Historical Stations}

\subsubsection{Origins: The World's First Broadcaster?}

Charles Herrold of San Jose was broadcasting music and talk on a regular schedule between 1912
and 1917, to a growing audience and to receiving stations at the 1915 Panama-Pacific International
Exhibition. His station, which became KCBS (740~AM), is one of the strongest candidates for the
title of the world's first continuously operating broadcast station---a claim that predates
KDKA Pittsburgh's November 1920 commercial debut. Following World War~I's radio silence, it
reemerged and eventually became a CBS-owned-and-operated news powerhouse.

In April 1920, the Western Radio Electric Company received a license for experimental station
6XD in Los Angeles (920~kHz), described by Pacific Radio News in 1921 as broadcasting on 325
meters two nights a week. This was the first broadcasting station in Southern California, active
approximately five months before KNX's predecessor, 6ADZ.

\subsubsection{Key California AM Stations: Profiles}

\begin{center}
\rowcolors{2}{rowgray}{white}
\begin{tabularx}{\textwidth}{L{2.5cm}C{1.5cm}L{2.5cm}X}
\toprule
\rowcolor{primary}
\tableheader{Station} & \tableheader{Freq.} & \tableheader{Founded} & \tableheader{Key History} \\
\midrule
KGO/KSFO SF & 810 & Jan 8, 1924 & GE/RCA/NBC Blue/ABC; Jack LaLanne; 50 kW Class A; call changed to KSFO in 2024 \\
KNBR SF & 680 & 1922 & NBC Red flagship SF; sports powerhouse; Giants/Warriors \\
KYA SF & 1260 & Dec 17, 1926 & Pacific Broadcasting Corp; CBS affiliate 1928 \\
KCBS SF & 740 & 1909* & World's first continuous broadcaster (Herrold); CBS O\&O; news \\
KFI Los Angeles & 640 & Apr 16, 1922 & First 50 kW west of Chicago (1931); Class A; Earle C. Anthony/NBC \\
KNX Los Angeles & 1070 & 1920 & Westinghouse/CBS O\&O; AP/UPI; Breaking News Murrow Award 2017 \\
KABC Los Angeles & 790 & 1925 & NBC Blue/ABC; news-talk; iHeartMedia \\
KMPC Los Angeles & 710 & 1922 & Music of Your Life; sports; later KSPN \\
\bottomrule
\end{tabularx}
\end{center}

\subsubsection{KGO/KSFO: The West Coast Clear-Channel Giant}

KGO signed on January 8, 1924, from General Electric's Oakland transformer plant as part of a
planned three-station GE network (with WGY Schenectady and KOA Denver). It was known as the
``Sunset Station.'' Due to GE's involvement with RCA, KGO was soon operated under NBC management.
The 1928 Federal Radio Commission reallocation designated clear-channel status on 810~kHz
originally at 790~kHz; NARBA in 1941 moved it to 810.

In 1943, the FCC forced NBC to sell one of its two networks; the Blue Network became ABC, and
KGO became ABC's West Coast flagship. In the postwar period, KGO produced live Western Swing
music and hosted the first radio exercise show by Jack LaLanne.

By the 1960s KGO had reinvented itself as one of America's first news-talk stations under Program
Director Jim Dunbar, who broadcast live from local restaurants and nightclubs. It held Arbitron's
\#1 Bay Area ranking for over 27 consecutive years. In 2024, the ``KGO'' call letters were retired
and the station became KSFO---a notable end to a century-long call sign history.

\subsubsection{KFI: The First 50,000-Watt Station West of Chicago}

KFI (640~AM) was first licensed March 31, 1922, by Earle~C.~Anthony, owner of a Packard
automobile dealership who had trained as an electrical engineer. It made its debut broadcast
on April 16, 1922, with vaudeville performers Eugene and Willie Howard.

In July 1931, KFI increased its transmitter power to 50,000 watts---the first station west of
Chicago to broadcast at that power. A special four-hour program marked the occasion, with
congratulatory speeches joined by NBC entertainers on a coast-to-coast hookup. KFI was part
of the NBC Pacific Orange Network alongside KGO, KGW, KOMO, and KHQ.

Today KFI (owned by iHeartMedia) remains one of the highest-rated talk stations on the West
Coast, with Bill Handel's morning show running since 1993. Its signal covers Southern California
by day and reaches much of western North America at night.

\subsubsection{The Don Lee Network and California Infrastructure}

Don Lee, a Cadillac dealer turned broadcaster, built the West Coast's most extensive broadcast
infrastructure in the 1930s. The Don Lee Network produced programming for the Mutual Broadcasting
System and served as its West Coast backbone. Don Lee stations included KFRC San Francisco,
KHJ Los Angeles, and affiliates throughout California and into the Pacific Northwest. Don Lee
became a stockholder in Mutual in 1940.

% ── MODULE 4: SPECTRUM AND LICENSING ─────────────────────────────────────────
\subsection{Module 4 Reference: Spectrum Availability and Licensing Pathways}

\subsubsection{The Spectrum Scarcity Reality}

The FCC states explicitly: ``In many areas of the country, no frequencies may be available on
which a new station could begin operating without causing interference to existing stations.''
This is especially acute in the top 50 markets---Seattle, Portland, San Francisco, Los Angeles,
and San Diego---where the FM dial is heavily occupied and AM frequencies are constrained by
Class~A clear-channel protections and NARBA international obligations.

The FM band (88.1--107.9~MHz) cannot be expanded: aeronautical navigation operations occupy
108--136~MHz above, and Channel~6 TV operations (82.0--88.0~MHz) constrain the lower edge.
The AM band was last expanded in the 1990s from 1600 to 1705~kHz, and further expansion is
described by the FCC as unlikely.

\subsubsection{Licensing Pathway Comparison}

\begin{center}
\rowcolors{2}{rowgray}{white}
\begin{tabularx}{\textwidth}{L{2.2cm}XL{2.5cm}L{2.2cm}}
\toprule
\rowcolor{primary}
\tableheader{Type} & \tableheader{Key Requirements} & \tableheader{Power / Range} & \tableheader{Availability} \\
\midrule
Full-Power AM & Citizenship, character, technical engineering study, construction permit, FCC Form~301 & 250 W minimum; Class A up to 50 kW & Very limited in major markets \\
Full-Power FM & Engineering study, open allotment, FCC Form~301, application fee & 100 W minimum (Class A) to 100 kW (Class C) & Very limited in major markets \\
LPFM & Nonprofit/educational entity; community-based (10-mile rule); FCC Form 318 during filing window & Up to 100 W; 3.5 mile radius & Filing windows \textasciitilde{}every 10 years; last window Nov--Dec 2023 \\
FM Translator & Rebroadcasts another station; FCC Form 349 & Typically 1--250 W & Limited; AM Revitalization opened access \\
Part 15 Unlicensed & No license required; must accept all interference & 200-foot effective radius & Always available; minimal coverage \\
\bottomrule
\end{tabularx}
\end{center}

\subsubsection{LPFM: The Community Broadcaster's Pathway}

The LPFM service was created by the FCC in January 2000. LPFM stations are authorized for
noncommercial educational broadcasting only (no commercial operation) and operate with up to
100 watts ERP at 30 meters antenna height, with an approximate service radius of 3.5 miles
(5.6~km). LPFM stations are not protected from interference from full-power stations.

Eligibility is restricted to nonprofits and governmental/educational entities that are
community-based (physically headquartered or 75\% of board members within 10 miles of the
proposed transmitter; 20 miles outside the top 50 markets). Individuals and commercial entities
are ineligible. Current broadcasters, cable operators, and newspaper publishers are also ineligible.

Applications are only accepted during FCC-announced filing windows. As of early 2026, the FCC
is NOT accepting new LPFM applications. The most recent window opened November 1, 2023, and
closed (revised) December 13, 2023---the third window ever, following 2000--2001 and 2013.
As of March 2023, there were 1,999 licensed LPFM stations in the United States. The next
window is expected approximately 10 years hence.

There is no filing fee for LPFM construction permit applications. Station construction and
operating costs vary; a basic LPFM setup can range from a few thousand to tens of thousands
of dollars depending on transmitter, antenna, tower, and studio requirements.

\subsubsection{Full-Power AM and FM: The Long Road}

For a commercial full-power AM station, applicants must: (1) identify an available frequency
using FCC AM Query tools; (2) commission a broadcast engineering study demonstrating no
interference to existing stations; (3) file FCC Form~301 (construction permit) electronically
via the Licensing and Management System (LMS) during an application window; (4) pay the
applicable filing fee; and (5) await FCC processing, which may involve comparative hearings
if competing applications are filed.

FCC licenses are issued for eight-year terms. The minimum AM power is 250 watts. Applications
must be filed electronically---paper forms are not accepted, though printed forms may be used
to prepare data. All applicants must be registered in the FCC's COmmission REgistration System
(CORES) with a valid FCC Registration Number (FRN).

\subsubsection{AM Revitalization: FM Translators for AM Stations}

As part of the ongoing AM Revitalization proceeding (Docket 13-249), the FCC has allowed AM
stations to obtain FM translator licenses to rebroadcast their programming on the FM band.
More than 2,100 FM translators were on the air rebroadcasting AM signals by 2020, with
approximately 700 more pending. This effectively allows daytime-only AM stations to maintain
a 24-hour presence on FM.

% ── MODULE 5: FCC RULES AND REGULATIONS ───────────────────────────────────────
\subsection{Module 5 Reference: FCC Rules and Regulations}

\subsubsection{The Regulatory Framework: 47 CFR Part 73}

All broadcast station rules are contained in 47 CFR Part~73, divided into subparts:
\begin{itemize}[leftmargin=*]
  \item \textbf{Subpart A} --- AM Broadcast Stations (\S\S~73.1--73.190)
  \item \textbf{Subpart B} --- FM Broadcast Stations (\S\S~73.201--73.333)
  \item \textbf{Subpart C} --- Digital Audio Broadcasting (\S\S~73.401--73.406)
  \item \textbf{Subpart D} --- Noncommercial Educational FM (\S\S~73.501--73.599)
  \item \textbf{Subpart G} --- Low Power FM Broadcast Stations (\S\S~73.801 et seq.)
  \item \textbf{Subpart H} --- Rules Applicable to All Broadcast Stations (\S\S~73.1001 et seq.)
\end{itemize}

Subpart~H is the licensee's primary compliance reference, covering political broadcasting,
EEO, EAS, indecency, and public file requirements.

\subsubsection{Public Interest Obligations}

The spectrum belongs to the public. Under the Radio Act of 1927 and the Communications Act
of 1934, licensees hold no property rights to the spectrum---only the privilege of using it
so long as they serve the public interest, convenience, and necessity. The FCC may revoke a
license before its term expires with notice and opportunity to be heard.

Key obligations for all broadcast licensees:
\begin{itemize}[leftmargin=*]
  \item \textbf{Equal Employment Opportunity (EEO):} Stations with five or more full-time employees
        must maintain outreach programs for employment and file FCC Form~396 (EEO Report) at
        license renewal.
  \item \textbf{Political Broadcasting:} During candidate elections, stations must provide equal
        time (``equal opportunity'') to legally qualified candidates for the same office under
        47 U.S.C.~\S~315; must offer lowest unit rate during campaign windows.
  \item \textbf{Emergency Alert System (EAS):} All AM and FM stations must participate in the EAS,
        which requires FCC-compliant EAS decoders/encoders. KFI and KNX serve as primary EAS
        entry points for the Southern California Emergency Alert System.
  \item \textbf{Indecency and Obscenity:} 18 U.S.C.~\S~1464 prohibits obscene, indecent, and
        profane broadcasts. Indecent content (material that depicts or describes sexual or
        excretory organs in a patently offensive manner) is prohibited between 6:00~a.m.\ and
        10:00~p.m. local time. Obscene content is prohibited at all times.
  \item \textbf{Public Inspection File:} Stations must maintain a publicly accessible online
        file (OPIF) through the FCC containing the license, applications, EEO records, political
        file, and other documents. LPFM stations are exempt from the local public file requirement
        but must maintain a political file.
  \item \textbf{On-Air Announcements:} During license renewal, stations must make on-air
        announcements about their renewal application within five business days of FCC public
        notice that the application has been accepted for filing.
\end{itemize}

\subsubsection{License Terms and Renewal}

FCC broadcast licenses are issued for eight-year terms. All radio station licenses are scheduled
to expire between 2027 and 2030, on a state-by-state schedule. Renewal applications (FCC
Form~303-S) must be filed four months before the expiration date. The FCC may not automatically
renew a license; the public may file Petitions to Deny or informal objections.

\subsubsection{Ownership Rules}

The Telecommunications Act of 1996 eliminated national radio station ownership caps. Local
ownership is governed by numerical market tiers based on total commercial stations in a market:
\begin{itemize}[leftmargin=*]
  \item In markets with 45+ stations: one entity may own up to 8 stations (no more than 5 in one service, AM or FM)
  \item In markets with 30--44 stations: up to 7 (4 per service)
  \item In markets with 15--29 stations: up to 6 (4 per service)
  \item In markets with 14 or fewer: up to 5 (3 per service)
\end{itemize}

The FCC's cross-ownership rules and the dual-network rule (\S~73.658) further constrain media
consolidation. The dual-network rule prohibits common ownership of two of the ``top four''
networks (ABC, CBS, Fox, NBC).

\subsubsection{Technical Standards}

AM stations must use FCC-approved transmitting equipment and must operate during hours
appropriate to their frequency class. Class~D (daytime) stations must cease or reduce power
at local sunset. Directional antenna arrays require FCC approval of the antenna system.
The minimum AM power for a construction permit is 250 watts.

The March 2025 FCC rulemaking (MB Docket 24-626) proposed updates to Part~73 rules to
reflect current application processing requirements, update form names, remove CDBS references
in favor of the LMS system, and eliminate certain AM power increase restrictions.

% ── MODULE 6: DIGITAL FUTURE ──────────────────────────────────────────────────
\subsection{Module 6 Reference: The Digital Future of AM Radio}

\subsubsection{The AM Revitalization Proceeding (Docket 13-249)}

In October 2013, the FCC issued a Notice of Proposed Rulemaking titled ``Revitalization of the
AM Radio Service.'' The proceeding has produced multiple Report and Orders, including:
\begin{itemize}[leftmargin=*]
  \item \textbf{First R\&O (2015):} Changes to daytime/nighttime community coverage standards;
        elimination of the AM ``ratchet'' rule; easier MDCL (Modulation Dependent Carrier Level)
        requirements; relaxation of minimum antenna efficiency standards.
  \item \textbf{Second R\&O:} FM translator access for AM stations.
  \item \textbf{Third R\&O (2017):} Relaxation of directional antenna array requirements.
  \item \textbf{Auction 99 (2018):} FM translator auction specifically for AM stations.
  \item \textbf{All-Digital R\&O (2020):} Authorized voluntary all-digital AM operation.
\end{itemize}

\subsubsection{HD Radio and All-Digital AM}

HD Radio (In-Band On-Channel, or IBOC digital) operates as a hybrid system, adding digital
sidebands to the existing analog AM carrier. The FM digital power increase order approved
by the FCC in 2010 and the asymmetric digital sideband order of September 2024 have expanded
HD Radio capabilities on FM. On AM, the HD Radio hybrid has faced adoption challenges.

In December 2020, the FCC approved voluntary all-digital AM operation using the MA3 mode
(the HD Radio standard for all-digital). This allows AM stations to cease analog transmission
and broadcast entirely in digital, improving audio quality and signal robustness. Hubbard
Broadcasting's WWFD in Frederick, Maryland, conducted the landmark experiments; testing showed
reliable daytime coverage to the 0.5~mV predicted contour and nighttime coverage to the NIF
contour. As of 2020, approximately 95 million vehicles on U.S. roads were equipped with HD
Radio receivers.

\subsubsection{The EV Interference Crisis}

Electric vehicles generate electromagnetic interference that degrades AM reception. Automakers
including Ford, BMW, and Volkswagen began removing AM radio from some EV models in 2023.
Ford initially announced removal from 2024 models; after consumer, broadcaster, and congressional
backlash, Ford CEO Jim Farley reversed the decision. The National Association of Broadcasters
urged all automakers to maintain AM capability. Congressional activity produced the AM Radio
for Every Vehicle Act, advancing through committee in 2023--2024.

GM and Stellantis announced no plans to remove AM from their vehicles, while keeping product
planning details confidential. The underlying technical problem---EV motor and charging
system emissions in the AM frequency range---remains unresolved.

\subsubsection{The Competitive Landscape}

As of March 2024, there were 17,370 licensed terrestrial stations in the United States (FCC,
Communications Marketplace Report 2024). AM/FM radio's combined share of total audio listening
fell to 36\% in 2023 (Edison Research), while online audio platforms gained ground. Terrestrial
radio revenue reached \$15.2 billion in 2023, down from \$17.8 billion in 2019.

In-car listening remains AM/FM radio's primary stronghold. Monthly podcast listenership rose
to 47\% of the U.S. population in 2024, with podcast ad revenue reaching \$1.9 billion in 2023.
Spotify leads audio streaming with 35\% of U.S. listeners aged 12+. SiriusXM had 33.9 million
subscribers in 2023 (\$6.8 billion revenue).

Despite these pressures, AM broadcasting serves critical roles that streaming cannot replicate:
emergency alert infrastructure, hyperlocal community service in remote areas, and minority-language
programming. In California alone, 34\% of AM stations serve specific minority communities (as
documented in the AM Revitalization proceeding comments).

\subsubsection{The Path Forward: Practical Options for AM Stations}

\begin{enumerate}[leftmargin=*, label=\textbf{\arabic*.}]
  \item \textbf{All-digital conversion:} File FCC Form~335-AM Digital Notification; commence
        all-digital operation 30 days after FCC public notice. Requires transmitter capable of MA3.
  \item \textbf{FM translator acquisition:} File for an FM translator to simulcast AM content on FM,
        providing 24/7 coverage even for daytimer stations.
  \item \textbf{HD Radio hybrid:} Deploy IBOC digital sidebands on existing AM analog authorization.
  \item \textbf{Internet simulcast:} Stream AM programming online at no additional FCC licensing cost.
  \item \textbf{Emergency infrastructure positioning:} Emphasize EAS role and community resilience
        value to justify ongoing public interest obligations.
\end{enumerate}

\subsection{Safety and Sensitivity Considerations}

\begin{center}
\rowcolors{2}{rowgray}{white}
\begin{tabularx}{\textwidth}{L{3.5cm}XC{2.5cm}}
\toprule
\rowcolor{primary}
\tableheader{Area} & \tableheader{Consideration} & \tableheader{Boundary Type} \\
\midrule
Source quality & All claims must cite FCC documents, peer-reviewed research, or professional broadcast organizations. No entertainment media, blogs, or unsourced claims. & ABSOLUTE \\
Numerical attribution & Every wattage, frequency, station count, or audience share attributed to specific source. & ABSOLUTE \\
No policy advocacy & Document presents regulatory information and findings without advocating for specific FCC policy outcomes. & ABSOLUTE \\
Legal advice boundary & Regulatory summaries are informational only; applicants must consult a broadcast attorney for specific licensing advice. & ANNOTATE \\
Minority programming & Statistics on minority-language and minority-community AM stations presented factually and without characterization. & GATE \\
\bottomrule
\end{tabularx}
\end{center}

\subsection{Source Bibliography}

\subsubsection{Government and Agency Sources}
\begin{itemize}[leftmargin=*]
  \item Federal Communications Commission. \textit{How to Apply for a Radio or Television Broadcast Station.} fcc.gov/media/radio/how-to-apply
  \item Federal Communications Commission. \textit{Low Power FM (LPFM) Broadcast Radio Stations.} fcc.gov/media/radio/lpfm
  \item Federal Communications Commission. \textit{AM Revitalization.} fcc.gov/media/radio/am-revitalization
  \item Federal Communications Commission. \textit{Digital Radio.} fcc.gov/media/radio/digital-radio
  \item Federal Communications Commission. \textit{AM Station Classes, and Clear, Regional, and Local Channels.} fcc.gov/media/radio/am-clear-regional-local-channels
  \item Federal Communications Commission. \textit{All-Digital AM Broadcasting, Revitalization of the AM Radio Service.} FCC 20-154 (Report and Order, Dec.\ 2020). Federal Register, Dec.\ 3, 2020.
  \item Federal Communications Commission. \textit{Broadcast Station Rule Updates.} MB Docket 24-626. Federal Register, Mar.\ 24, 2025.
  \item eCFR. \textit{47 CFR Part 73 --- Radio Broadcast Services.} ecfr.gov/current/title-47/chapter-I/subchapter-C/part-73
  \item FCC. \textit{Communications Marketplace Report 2024.} (Cited in Radio Ink, Jan.\ 7, 2025)
\end{itemize}

\subsubsection{Professional Organizations and Research}
\begin{itemize}[leftmargin=*]
  \item California Historical Radio Society. \textit{100 Years of Radio.} californiahistoricalradio.com (2011, updated 2024)
  \item HistoryLink.org. Banel, Feliks. \textit{Radio in Washington.} historylink.org/File/22494
  \item Bay Area Radio Museum \& Hall of Fame. \textit{KYA/KOIT Historical Timeline.} bayarearadio.org
  \item Radio Survivor. \textit{Radio Recollections: The San Francisco Bay Area in the 1950s.} radiosurvivor.com (June 2016)
  \item Common Frequency. \textit{Overview on Applying for Low Power FM in 2023.} commonfrequency.org
  \item Xperi. \textit{Is AM Radio Dying? Not Any Time Soon.} xperi.com/blog (2023)
  \item Radio Ink. \textit{FCC Captures Radio's Struggle Against Digital in Biennial Report.} radioink.com (Jan.\ 7, 2025)
  \item Radio World. \textit{FCC Approves All-Digital Option for AM.} radioworld.com (Oct.\ 27, 2020)
\end{itemize}

\subsubsection{Station Histories and Archives}
\begin{itemize}[leftmargin=*]
  \item Wikipedia contributors. \textit{KFI.} en.wikipedia.org/wiki/KFI (updated Feb.\ 2026)
  \item Wikipedia contributors. \textit{KSFO.} en.wikipedia.org/wiki/KSFO (updated 2026)
  \item Wikipedia contributors. \textit{KNWN (AM).} en.wikipedia.org/wiki/KNWN\_(AM) (updated Jan.\ 2026)
  \item QZVX.com. \textit{AM Radio Stations: PNW History.} qzvx.com/am-radio-stations/
  \item Swingin' West. \textit{Oldest Radio Stations / American / Radio History.} swinginwest.com
  \item Tripod / Jeff560. \textit{A Chronology of AM Radio Broadcasting 1900--1960.} jeff560.tripod.com/chrono1.html
  \item Oregon Radio / KGW / NBC History. ke7ggv.tripod.com/nwradiohistory.html
\end{itemize}

\newpage

% ══════════════════════════════════════════════════════════════════════════════
% STAGE 3 — MISSION PACKAGE
% ══════════════════════════════════════════════════════════════════════════════
\stageheader{STAGE 3 --- MISSION PACKAGE}{tertiary}

\section{Milestone Specification}

\subsection{Mission Objective}

Produce a comprehensive, publication-quality research document on the history, present state,
and future of AM radio broadcasting on the West Coast of the United States, with deep focus on
the Pacific Northwest and California. The document will serve both as an authoritative historical
reference and as a practical licensing and regulatory guide for prospective broadcasters in 2026.
All claims will be sourced to FCC documents, professional broadcast organizations, or peer-reviewed
research.

\subsection{Architecture Overview}

\begin{asciibox}
WAVE EXECUTION ARCHITECTURE
══════════════════════════════════════════════════════════════════

  WAVE 0: FOUNDATION
  ┌──────────────────────────────────────────────────────────┐
  │  W0-A: Shared Schema (Haiku)   W0-B: Source Index (Sonnet)│
  └──────────────────────────────────────────────────────────┘
                              │
  WAVE 1: PARALLEL RESEARCH TRACKS
  ┌──────────────────────────────────────────────────────────┐
  │  TRACK ALPHA                 │  TRACK BETA               │
  │  W1-A: AM Band Foundations   │  W1-D: Spectrum & Licensing│
  │  W1-B: PNW Station History   │  W1-E: FCC Regulations    │
  │  W1-C: California Stations   │                           │
  └──────────────────────────────────────────────────────────┘
                              │
  WAVE 2: SYNTHESIS
  ┌──────────────────────────────────────────────────────────┐
  │  W2-A: Cross-Module Integration (Opus)                   │
  │  W2-B: Digital Future Analysis (Opus)                    │
  └──────────────────────────────────────────────────────────┘
                              │
  WAVE 3: PUBLICATION + VERIFICATION
  ┌──────────────────────────────────────────────────────────┐
  │  W3-A: Document Assembly (Sonnet)                        │
  │  W3-B: Source Verification (Sonnet)                      │
  │  W3-C: CAPCOM Review Gate (Human)                        │
  └──────────────────────────────────────────────────────────┘
                              │
                    FINAL DELIVERABLE PDF
\end{asciibox}

\subsection{System Layers}

\begin{enumerate}[leftmargin=*, label=\textbf{\arabic*.}]
  \item \textbf{Foundation Layer} --- Shared data schemas, source index, station database structure
  \item \textbf{Historical Research Layer} --- Station-by-station profiles, network chronologies, format evolution
  \item \textbf{Regulatory Research Layer} --- FCC rule summaries, licensing pathways, spectrum analysis
  \item \textbf{Synthesis Layer} --- Cross-module integration, digital future analysis, practical broadcaster guide
  \item \textbf{Publication Layer} --- Document assembly, formatting, source verification, final PDF
\end{enumerate}

\subsection{Deliverables}

\begin{center}
\rowcolors{2}{rowgray}{white}
\begin{tabularx}{\textwidth}{C{0.5cm}L{3.5cm}XL{3cm}}
\toprule
\rowcolor{primary}
\tableheader{\#} & \tableheader{Deliverable} & \tableheader{Acceptance Criteria} & \tableheader{Primary Spec} \\
\midrule
1 & Station History Compendium & All Class A/B PNW and CA stations profiled; dates, frequencies, affiliations verified & W1-B, W1-C \\
2 & AM Band Technical Reference & Band structure, propagation physics, station classes, NARBA explained & W1-A \\
3 & Licensing Pathway Guide & All five pathways documented; step-by-step; current as of 2026 & W1-D \\
4 & FCC Regulatory Summary & 47 CFR Part 73 obligations for new licensees; ownership rules; EAS & W1-E \\
5 & Digital Transition Analysis & HD Radio, all-digital AM, FM translators, EV interference, market data & W2-B \\
6 & Practical Broadcaster Road Map & Synthesized entry guide for prospective 2026 West Coast broadcaster & W2-A \\
7 & Verified Source Bibliography & All sources traceable; no entertainment media & W3-B \\
\bottomrule
\end{tabularx}
\end{center}

\subsection{Component Breakdown}

\begin{center}
\rowcolors{2}{rowgray}{white}
\begin{tabularx}{\textwidth}{L{2.5cm}L{2.5cm}C{1.5cm}C{2cm}}
\toprule
\rowcolor{primary}
\tableheader{Component} & \tableheader{Dependencies} & \tableheader{Model} & \tableheader{Token Budget (K)} \\
\midrule
W0-A: Schema & None & Haiku & 4 \\
W0-B: Source Index & None & Sonnet & 8 \\
W1-A: AM Foundations & W0 & Sonnet & 18 \\
W1-B: PNW Stations & W0, W1-A & Sonnet & 28 \\
W1-C: CA Stations & W0, W1-A & Sonnet & 28 \\
W1-D: Licensing & W0, W1-A & Sonnet & 22 \\
W1-E: FCC Rules & W0, W1-A & Sonnet & 22 \\
W2-A: Cross-Module Integration & W1-B, W1-C, W1-D, W1-E & Opus & 35 \\
W2-B: Digital Future & W1-A, W2-A & Opus & 25 \\
W3-A: Document Assembly & W2-A, W2-B & Sonnet & 40 \\
W3-B: Source Verification & W3-A & Sonnet & 16 \\
\textbf{Total} & & & \textbf{246} \\
\bottomrule
\end{tabularx}
\end{center}

\subsection{Model Assignment Rationale}

\begin{center}
\rowcolors{2}{rowgray}{white}
\begin{tabularx}{\textwidth}{L{2.5cm}C{1.5cm}X}
\toprule
\rowcolor{primary}
\tableheader{Assignment} & \tableheader{Model} & \tableheader{Rationale} \\
\midrule
Schema, scaffolding & Haiku & Mechanical, low-judgment template generation \\
Survey, compilation, verification & Sonnet & Standard research and document production \\
Cross-module synthesis, digital future analysis, broadcaster road map & Opus & Requires judgment, synthesis across 5+ modules, trend analysis and extrapolation \\
\bottomrule
\end{tabularx}
\end{center}

\textbf{Model allocation:} Opus $\approx$24\% (60K tokens), Sonnet $\approx$72\% (178K tokens), Haiku $\approx$4\% (8K tokens).

\section{Wave Execution Plan}

\subsection{Wave Summary}

\begin{center}
\rowcolors{2}{rowgray}{white}
\begin{tabularx}{\textwidth}{C{1cm}L{3cm}XC{2cm}C{2cm}}
\toprule
\rowcolor{primary}
\tableheader{Wave} & \tableheader{Name} & \tableheader{Components} & \tableheader{Mode} & \tableheader{Gate} \\
\midrule
0 & Foundation & W0-A, W0-B & Sequential & Auto \\
1 & Parallel Research & W1-A through W1-E & Parallel (2 tracks) & CAPCOM \\
2 & Synthesis & W2-A, W2-B & Sequential & CAPCOM \\
3 & Publication & W3-A, W3-B, W3-C & Sequential & Human \\
\bottomrule
\end{tabularx}
\end{center}

\subsection{Wave 0: Foundation}

\begin{center}
\rowcolors{2}{rowgray}{white}
\begin{tabularx}{\textwidth}{L{2cm}L{3cm}XC{2cm}}
\toprule
\rowcolor{primary}
\tableheader{Task} & \tableheader{Agent} & \tableheader{Deliverable} & \tableheader{Model} \\
\midrule
W0-A & PLAN + Haiku & Station database schema (JSON): call sign, frequency, city, class, founded, network affiliations, format history, current status & Haiku \\
W0-B & SCOUT + Sonnet & Annotated source index: URL, document type, reliability tier, module assignments for each of the 30+ sources gathered in cold-start research & Sonnet \\
\bottomrule
\end{tabularx}
\end{center}

\subsection{Wave 1: Parallel Research Tracks}

\textbf{Track Alpha --- History}

\begin{center}
\rowcolors{2}{rowgray}{white}
\begin{tabularx}{\textwidth}{L{2cm}L{2.5cm}XC{2cm}}
\toprule
\rowcolor{primary}
\tableheader{Task} & \tableheader{Agent} & \tableheader{Deliverable} & \tableheader{Model} \\
\midrule
W1-A & CRAFT-TECH & AM band foundations document: physics, FRC/FCC history, NARBA, station classification table, band expansion history & Sonnet \\
W1-B & CRAFT-HIST + EXEC-1 & PNW station compendium: full profiles for all Class A/B stations in WA/OR/ID; network chronology; format evolution by decade & Sonnet \\
W1-C & CRAFT-HIST + EXEC-2 & California station compendium: full profiles for Bay Area and LA Class A/B stations; Don Lee Network; NBC Pacific structures & Sonnet \\
\bottomrule
\end{tabularx}
\end{center}

\textbf{Track Beta --- Regulatory \& Licensing}

\begin{center}
\rowcolors{2}{rowgray}{white}
\begin{tabularx}{\textwidth}{L{2cm}L{2.5cm}XC{2cm}}
\toprule
\rowcolor{primary}
\tableheader{Task} & \tableheader{Agent} & \tableheader{Deliverable} & \tableheader{Model} \\
\midrule
W1-D & CRAFT-REG & Licensing pathway guide: all five pathways with step-by-step, forms, fees, LPFM window schedule, spectrum scarcity context for 4 major West Coast markets & Sonnet \\
W1-E & CRAFT-REG & FCC regulatory summary: Part 73 subpart-by-subpart; public interest obligations; ownership rules; EAS; indecency; license terms; 2025 rule updates & Sonnet \\
\bottomrule
\end{tabularx}
\end{center}

\subsection{Wave 2: Synthesis}

\begin{center}
\rowcolors{2}{rowgray}{white}
\begin{tabularx}{\textwidth}{L{2cm}L{2.5cm}XC{2cm}}
\toprule
\rowcolor{primary}
\tableheader{Task} & \tableheader{Agent} & \tableheader{Deliverable} & \tableheader{Model} \\
\midrule
W2-A & INTEG + Opus & Cross-module integration: consistent station data; practical broadcaster road map synthesizing history + licensing + regulations; editorial through-line & Opus \\
W2-B & ANALYST + Opus & Digital future analysis: HD Radio adoption trajectory; all-digital AM viability; EV interference technical analysis; competitive streaming landscape; 5-year AM outlook & Opus \\
\bottomrule
\end{tabularx}
\end{center}

\subsection{Wave 3: Publication and Verification}

\begin{center}
\rowcolors{2}{rowgray}{white}
\begin{tabularx}{\textwidth}{L{2cm}L{2.5cm}XC{2cm}}
\toprule
\rowcolor{primary}
\tableheader{Task} & \tableheader{Agent} & \tableheader{Deliverable} & \tableheader{Model} \\
\midrule
W3-A & EXEC-1 + Sonnet & Full document assembly: LaTeX or structured markdown of all modules; tables, diagrams, bibliography formatted to publication standard & Sonnet \\
W3-B & VERIFY + Sonnet & Source verification report: every factual claim traced to source; SC-SRC and SC-NUM safety tests run; flagged items returned to producing agent & Sonnet \\
W3-C & CAPCOM & Human review gate: final HITL approval before PDF generation & Human \\
\bottomrule
\end{tabularx}
\end{center}

\subsection{Cache Optimization Strategy}

\begin{itemize}[leftmargin=*]
  \item Cache the annotated source index (W0-B) at session start for all Wave~1 agents---avoids 30+ URL lookups per agent
  \item Cache the AM band foundations document (W1-A) as shared context for all history and regulatory agents
  \item Maintain the station database schema (W0-A) in shared context throughout Wave~1 to ensure consistent data structure
  \item Batch all FCC LMS queries within a single SCOUT agent session in Wave~0 rather than distributing across agents
  \item Use prompt caching for the 47 CFR Part~73 regulatory text---accessed by both W1-D and W1-E
\end{itemize}

\subsection{Token Budget}

\begin{center}
\rowcolors{2}{rowgray}{white}
\begin{tabularx}{\textwidth}{L{3cm}C{2cm}C{2cm}C{2cm}C{3cm}}
\toprule
\rowcolor{primary}
\tableheader{Wave} & \tableheader{Opus (K)} & \tableheader{Sonnet (K)} & \tableheader{Haiku (K)} & \tableheader{Wave Total (K)} \\
\midrule
Wave 0 & 0 & 8 & 4 & 12 \\
Wave 1 (Alpha) & 0 & 74 & 0 & 74 \\
Wave 1 (Beta) & 0 & 44 & 0 & 44 \\
Wave 2 & 60 & 0 & 0 & 60 \\
Wave 3 & 0 & 56 & 0 & 56 \\
\textbf{Total} & \textbf{60} & \textbf{182} & \textbf{4} & \textbf{246} \\
\bottomrule
\end{tabularx}
\end{center}

\subsection{Dependency Graph}

\begin{asciibox}
DEPENDENCY GRAPH
══════════════════════════════════════════════════════

  [W0-A Schema] ─────────────────────────────────────┐
  [W0-B Source Index] ──┬───────────────────────────────┤
                        │                               │
            ┌───────────┼──────────┬──────────┐        │
            ▼           ▼          ▼          ▼        ▼
         [W1-A]     [W1-B]     [W1-C]    [W1-D]   [W1-E]
         Foundations PNW Hist  CA Hist  Licensing  FCC Rules
            │           │          │          │        │
            └───────────┴──────────┴──────────┴────────┘
                                  │
                    ┌─────────────┴──────────────┐
                    ▼                            ▼
                 [W2-A]                       [W2-B]
                Cross-Module               Digital Future
                Integration                  Analysis
                    │                            │
                    └─────────────┬──────────────┘
                                  │
                    ┌─────────────┴──────────────┐
                    ▼                            ▼
                 [W3-A]                       [W3-B]
                Document                     Source
                Assembly                  Verification
                    │                            │
                    └─────────────┬──────────────┘
                                  │
                              [W3-C]
                           CAPCOM GATE
                                  │
                        FINAL DELIVERABLE
\end{asciibox}

\section{Test Plan}

\subsection{Test Categories}

\begin{center}
\rowcolors{2}{rowgray}{white}
\begin{tabularx}{\textwidth}{L{3cm}C{2cm}L{2cm}C{2cm}L{3cm}}
\toprule
\rowcolor{primary}
\tableheader{Category} & \tableheader{Count} & \tableheader{Priority} & \tableheader{Action} & \tableheader{Target \%} \\
\midrule
Safety-critical & 7 & Mandatory & BLOCK & 16\% \\
Core functionality & 22 & Required & BLOCK & 49\% \\
Integration & 10 & Required & BLOCK & 22\% \\
Edge cases & 6 & Best-effort & LOG & 13\% \\
\textbf{Total} & \textbf{45} & & & \textbf{100\%} \\
\bottomrule
\end{tabularx}
\end{center}

\subsection{Safety-Critical Tests}

\begin{center}
\rowcolors{2}{rowgray}{white}
\begin{tabularx}{\textwidth}{C{1.5cm}L{3.5cm}X}
\toprule
\rowcolor{primary}
\tableheader{Test ID} & \tableheader{Verifies} & \tableheader{Expected} \\
\midrule
SC-SRC & Source quality & All citations traceable to FCC documents, peer-reviewed sources, or professional broadcast organizations. Zero entertainment media or unsourced blogs. \\
SC-NUM & Numerical attribution & Every wattage, frequency, station count, audience share, and date attributed to a specific named source. \\
SC-ADV & No policy advocacy & Document presents regulatory information without advocating for specific FCC policy changes or positions. \\
SC-LEG & Legal advice boundary & All regulatory summaries include disclaimer that content is informational and applicants must consult a broadcast attorney. \\
SC-MIN & Minority community data & Statistics on minority-language/community AM stations presented factually; no characterization or editorial framing. \\
SC-CUR & Temporal currency & Document correctly reflects 2026 regulatory status; no outdated CDBS references; LMS system referenced for filings. \\
SC-ELI & LPFM eligibility accuracy & LPFM eligibility restrictions (nonprofits only; no individuals; no commercial entities) stated accurately with CFR citation. \\
\bottomrule
\end{tabularx}
\end{center}

\subsection{Core Functionality Tests (Selected)}

\begin{center}
\rowcolors{2}{rowgray}{white}
\begin{tabularx}{\textwidth}{C{1.5cm}L{3.5cm}X}
\toprule
\rowcolor{primary}
\tableheader{Test ID} & \tableheader{Verifies} & \tableheader{Expected} \\
\midrule
CF-01 & PNW station coverage completeness & All Class A and major Class B AM stations in WA/OR profiled with founding date, frequency, and ownership history \\
CF-02 & California station completeness & All Class A AM stations in CA (KFI, KNX, KGO/KSFO, KNBR, KCBS, KABC) individually profiled \\
CF-03 & Network chronology accuracy & NBC Pacific networks (Orange, Gold), Don Lee/Mutual, CBS, and ABC West Coast structures mapped with station affiliations and dates \\
CF-04 & AM band structure documentation & 535--1705~kHz range, 10~kHz channel spacing, Class A--D power limits, and daytime/nighttime propagation differences explained \\
CF-05 & LPFM pathway accuracy & Filing window schedule, Form 318, eligibility requirements, 100~W power limit, 3.5-mile radius, and community-based requirement all documented \\
CF-06 & Full-power AM pathway & Form 301, engineering study requirement, 250~W minimum, construction permit process, and FRN registration requirement documented \\
CF-07 & FM translator pathway & AM Revitalization translator access, 2,100+ on-air translators statistic sourced to FCC, Form 349 referenced \\
CF-08 & Part 15 unlicensed operation & 200-foot effective radius, FCC Part 15 rules, absence of licensing requirement, and interference acceptance noted \\
CF-09 & License term and renewal & Eight-year term, Form 303-S, four-month pre-filing deadline, 2027--2030 expiration schedule documented \\
CF-10 & Ownership rules & 1996 Telecom Act changes, local numerical market tiers, cross-ownership rules referenced \\
CF-11 & EAS obligations & EAS decoder/encoder requirement, KFI/KNX roles documented, CFR citation included \\
CF-12 & Indecency rules & 18 U.S.C.\ \S~1464, 6~a.m.--10~p.m.\ safe harbor, obscenity prohibition, and 47 CFR \S~73.3999 all referenced \\
CF-13 & All-digital AM authorization & FCC 20-154, MA3 mode, Form 335-AM Digital Notification, 30-day public notice window documented \\
CF-14 & HD Radio adoption data & 95 million vehicles with HD Radio sourced to Xperi/DTS 2023 data \\
CF-15 & EV interference documented & Automaker removal decisions (Ford, BMW, VW), Ford reversal, congressional response, NAB position documented \\
CF-16 & Audience share data & 36\% AM/FM share of audio listening in 2023 sourced to Edison Research via FCC 2024 report \\
CF-17 & AM revitalization history & Docket 13-249 milestones (2013 NPRM through 2020 digital R\&O) documented with dates \\
CF-18 & NARBA frequency shifts & KOMO 950$\to$1000, KGO 790$\to$810 and other NARBA-driven reassignments documented \\
CF-19 & KJR founding history & 1921 founding by Vincent Kraft, first PNW license, KYA San Francisco connection documented \\
CF-20 & KFI 50 kW milestone & July 1931 power increase, first west of Chicago, NBC Orange network affiliation documented \\
CF-21 & KGO/KSFO history & January 8, 1924 sign-on; GE/NBC/ABC lineage; Jack LaLanne first radio exercise show; KSFO call change 2024 \\
CF-22 & Spectrum expansion limitation & FM band constraints (aeronautical above, CH6 TV below), AM expanded band history documented \\
\bottomrule
\end{tabularx}
\end{center}

\subsection{Verification Matrix}

\begin{center}
\rowcolors{2}{rowgray}{white}
\begin{tabularx}{\textwidth}{XL{3.5cm}C{1.8cm}}
\toprule
\rowcolor{primary}
\tableheader{Success Criterion} & \tableheader{Test ID(s)} & \tableheader{Status} \\
\midrule
SC-1. All Class A/B PNW and CA stations profiled & CF-01, CF-02, CF-03 & Pending \\
SC-2. AM band structure fully explained & CF-04, CF-22 & Pending \\
SC-3. FCC licensing pathways documented & CF-05, CF-06, CF-07, CF-08 & Pending \\
SC-4. Spectrum scarcity context documented & CF-22, IN-01 & Pending \\
SC-5. 47 CFR Part 73 obligations summarized & CF-09, CF-11, CF-12 & Pending \\
SC-6. Ownership restrictions documented & CF-10 & Pending \\
SC-7. AM digital transition options described & CF-13, CF-14 & Pending \\
SC-8. EV interference documented & CF-15 & Pending \\
SC-9. Don Lee / NBC Pacific networks mapped & CF-03, CF-21 & Pending \\
SC-10. Historical format evolution documented & CF-17, CF-18, CF-19, CF-20 & Pending \\
SC-11. All numerical claims attributed & SC-NUM, CF-16 & Pending \\
SC-12. Practical broadcaster road map synthesized & IN-02, W2-A & Pending \\
\bottomrule
\end{tabularx}
\end{center}

\section{Mission Crew Manifest}

\begin{center}
\rowcolors{2}{rowgray}{white}
\begin{tabularx}{\textwidth}{L{2cm}L{3cm}X}
\toprule
\rowcolor{primary}
\tableheader{Role} & \tableheader{Agent} & \tableheader{Responsibility} \\
\midrule
FLIGHT & Orchestrator & Overall mission direction; go/no-go at each wave boundary \\
PLAN & Planner & Wave decomposition; parallel track optimization; dependency management \\
SCOUT & Research Agent & Source gathering; FCC database queries; web research; evidence compilation \\
EXEC-1 & Executor Alpha & Document production for Track Alpha (PNW history, California history) \\
EXEC-2 & Executor Beta & Document production for Track Beta (licensing, FCC rules) \\
CRAFT-HIST & History Specialist & Station-by-station chronicle; network affiliate research; archival sourcing \\
CRAFT-TECH & Technical Specialist & AM band physics; antenna systems; propagation; HD Radio technical standards \\
CRAFT-REG & Regulatory Specialist & 47 CFR Part 73 analysis; FCC LMS procedures; ownership rules; EAS \\
VERIFY & Verification Agent & Source validation; SC-SRC and SC-NUM safety tests; accuracy checking \\
INTEG & Integration Agent & Cross-module consistency; broadcaster road map synthesis; data reconciliation \\
ANALYST & Pattern Analyst & Digital future trend analysis; competitive landscape; 5-year outlook \\
CAPCOM & HITL Interface & Human approval at Wave 1 and Wave 2 boundaries; only H-M channel \\
BUDGET & Resource Manager & Token budget tracking; wave completion reporting; model allocation monitoring \\
SURGEON & Health Monitor & Research quality degradation detection; factual consistency monitoring \\
TOPO & Topology Manager & Agent team structure; mid-mission restructuring if needed \\
ARCH & Architecture Agent & Cross-cutting structural decisions; document template enforcement \\
LOG & Chronicle Agent & Commit messages; milestone summaries; audit trail \\
\bottomrule
\end{tabularx}
\end{center}

\section{Execution Summary}

\begin{center}
\rowcolors{2}{rowgray}{white}
\begin{tabularx}{\textwidth}{L{4cm}X}
\toprule
\rowcolor{primary}
\tableheader{Metric} & \tableheader{Value} \\
\midrule
Activation Profile & Fleet (6 modules, 17 roles) \\
Total Estimated Tokens & 246,000 \\
Model Split & Opus 24\% / Sonnet 72\% / Haiku 4\% \\
Wave Count & 4 waves \\
Parallel Efficiency & 2 simultaneous tracks in Wave 1 (72\% efficiency vs.\ sequential) \\
Estimated Sessions & 4--6 Claude Code sessions \\
Estimated Context Windows & 18--24 \\
Human Gate Points & 3 (Wave 1 start, Wave 2 start, final delivery) \\
Safety-Critical Tests & 7 (all mandatory-pass / BLOCK) \\
Total Tests & 45 \\
Primary Regulatory Authority & FCC / 47 CFR Part 73 \\
Research Domain & Broadcast history, regulatory law, spectrum engineering \\
Geographic Focus & Washington, Oregon, Idaho, California (West Coast primary) \\
Temporal Scope & 1909--2026 \\
\bottomrule
\end{tabularx}
\end{center}

\vspace{2em}

\begin{quotebox}
\textit{``Before the internet, the radio dial was the only place where a single voice could
reach the entire Pacific Northwest at once---in the dark, in the mountains, across the
water. The frequency was free. The license was granted by the public. The obligation was
to serve. That architecture still exists. The builder who learns to navigate it inherits
something older and stranger than any streaming platform: the right to speak to a whole
geography, and the responsibility to make that speech worth hearing.''}

\vspace{0.5em}
\hfill--- GSD Research Mission Vision, 2026 | Fox Infrastructure Group
\end{quotebox}

\end{document}
